\documentclass[english,10pt,a4paper]{article}

\usepackage{geometry}
\geometry{top=3cm, bottom=3cm, left=3cm , right=3cm}

\usepackage[utf8]{inputenc}
\usepackage[english]{babel}
\usepackage{hyperref}

\usepackage{amsmath,amssymb}
\usepackage{graphicx}
\usepackage{booktabs}
\usepackage{listings}

\usepackage{indentfirst}

\hypersetup{
    colorlinks=true,
    linkcolor=blue,
    filecolor=magenta,      
    urlcolor=cyan,
}



\begin{document}
\title{Statistical Characterisation of High-Field Magnet Operations}
\author{Valentin Levchenko \\
        Student ID: 22422485 \\
        \texttt{valentin.levchenko@etu.unistra.fr} \\ [0.5cm]
        \textbf{Supervisors:} \\
        Dr. Christophe Trophime (LNCMI, CNRS) \\
        Prof. Christophe Prud'homme (University of Strasbourg) \\
        Asst. Prof. Joubine Aghili (IRMA, University of Strasbourg)
}
\date{
    Master CSMI \\
    University of Strasbourg \\
    In collaboration with CNRS and LNCMI \\
    \today
}

\maketitle



\begin{abstract}


    GitHub repository:
    \url{https://github.com/master-csmi/2026-m1-hifimagnet.git}

    LNCMI website: \url{https://lncmi.cnrs.fr}

    EMFL website: \url{https://emfl.eu}
\end{abstract}



\newpage
\tableofcontents
\newpage



\section{Current progress}

\begin{itemize}
    \item Implemented autmoatic import of all 
          \texttt{hosuing-summary-YYYY.json} files (M9/M10, 2024-2026) into 
          DuckDB.
    \item Created \texttt{housing_summary} table and enriched it with \texttt{site} 
          and \texttt{year} metadata.
    \item Performed an initial data audit:
        \begin{itemize}
            \item 855 records imported
            \item no duplicate filenames detected
            \item identified missing archive, Pupitre and trigger file references
            \item verified tha t371 records can already be linked to existing
                  experiments throught the corresponding Pupitre file name
        \end{itemize}
    \item Verified and extrated the 2023 Pupitre files for future processing 
          and testing
    \item Set up the development environment to support the new record workflow
\end{itemize}


\pagebreak

\bibliographystyle{unsrt}
\bibliography{references}


\appendix

\section{User Guide}

The repository is organised into several main components:
\begin{itemize}
    \item \texttt{Dockerfile}
    \item \texttt{README.md}
    \item \texttt{requirements.txt}
    \item \texttt{build.md}
    \item \texttt{data}
    \item \texttt{docs}
    \item \texttt{results}
    \item \texttt{src}
\end{itemize}

The main analysis is implemented in the notebook \texttt{src/main.ipynb}.

Generated outputs are automatically exported into \texttt{results/<MAGNET\_NAME>/}.

The project was developed using a standard Python scientific environment including
\texttt{numpy, pandas, matplotlib} as well as additional libraries \texttt{duckdb, 
json, polars, great-tables}.

\subsection{Environment Setup}

Clone the repository:
\begin{verbatim}
    git clone https://github.com/master-csmi/2026-m1-hifimagnet.git
    cd 2026-m1-hifimagnet
\end{verbatim}

Create and activate a Python virtual environment:
\begin{verbatim}
    python3 -m venv .venv
    source .venv/bin/activate
\end{verbatim}

Install the required dependencies:
\begin{verbatim}
    pip install -r project1/requirements.txt
\end{verbatim}

Run the main notebook:
\begin{verbatim}
    jupyter notebook project1/src/main.ipynb
\end{verbatim}

\subsection{Docker}

Build the Docker image:
\begin{verbatim}
    docker build -t hifimagnet1 project1
\end{verbatim}

Run the container:
\begin{verbatim}
    docker run -it hifimagnet1
\end{verbatim}

\end{document}
